%% kapa: exam=1 round=28 year=2017 period=1 form=A subject=civil_law subject_ko=민법 seq=6 number=6
%% source: Q-Net 감정평가사 기출문제 (cst00309) artlSeq=5199603 file=extracted/1008/5199603/2017년 제28회 감정평가사 1차시험 문제지/1교시 A형.pdf
%% conversion: kordoc 4.24.51 -> tools/gen1.py (auto); review=proofread-sonnet; flags=none
\begin{question}{6}
비법인사단에 관한 설명으로 옳지 않은 것은? (다툼이 있으면 판례에 따름)
\begin{choices}
  \item 사단법인의 하부조직이라도 스스로 단체로서의 실체를 갖추고 독자적인 활동을 하고 있다면 그 사단법인과는 별개의 독립된 비법인사단으로 볼 수 있다.
  \item 정관 기타 규약에 다른 정함이 없는 한, 사원총회의 결의를 거치지 않은 총유물의 관리행위는 무효이다.
  \item 비법인사단의 대표자의 행위가 외관상ㆍ객관적으로 직무에 관한 행위로 인정될 수 있으면, 그의 행위가 직무에 관한 것이 아님을 피해자가 중대한 과실로 알지 못한 경우에도 비법인사단에게 손해배상책임을 물을 수 있다.
  \item 소집절차에 하자가 있어 그 효력을 인정할 수 없는 종중총회의 결의라도 후에 적법하게 소집된 종중총회에서 이를 추인하면 처음부터 유효로 된다.
  \item 재건축조합의 총회에서는 정관에 다른 정함이 없는 한 소집 1주간 전에 통지된 그 회의의 목적 사항에 관하여만 결의할 수 있다.
\end{choices}
\end{question}
\answer{3}
